\documentclass{article}
\usepackage{amsmath, amssymb}
\usepackage{geometry}
\geometry{margin=1in}

\begin{document}

\section*{Cybenko's Theorem (Universal Approximation Theorem)}

Let:
\[
K = [0,1]^n
\]

Define a single-hidden-layer network as follows:
\[
N(x) = \sum_{j=1}^{m} \alpha_j \, \sigma(w_j^T x + b_j)
\]
where:
\begin{itemize}
    \item \( x \in \mathbb{R}^n \)
    \item \( w_j \in \mathbb{R}^n \)
    \item \( b_j \in \mathbb{R} \)
    \item \( \alpha_j \in \mathbb{R} \)
    \item \( \sigma : \mathbb{R} \to \mathbb{R} \) is the activation function.
\end{itemize}

For the sigmoid function:
\[
\sigma(t) = \frac{1}{1 + e^{-t}}
\]

Cybenko's Theorem states that if \( \sigma \) is continuous and sigmoidal, meaning:
\[
\lim_{t \to -\infty} \sigma(t) = 0
\]
and
\[
\lim_{t \to +\infty} \sigma(t) = 1
\]
then the set
\[
\mathcal{N}_\sigma = \left\{ \sum_{j=1}^{m} \alpha_j \, \sigma(w_j^T x + b_j) \; : \; m < \infty \right\}
\]
is dense in the space
\[
C([0,1]^n)
\]
That is:
\[
\overline{\mathcal{N}_\sigma} = C([0,1]^n)
\]
And this is precisely the heart of the theorem.

\section*{What Does ``Dense'' Mean?}

This part is very important.

When we say
\[
\overline{\mathcal{N}_\sigma} = C(K)
\]
we do not mean that the network is exactly equal to the function.

Rather, for every
\[
f \in C(K)
\]
and every
\[
\varepsilon > 0
\]
one can find a network such that:
\[
\| f - N \|_\infty < \varepsilon
\]
where:
\[
\| f - N \|_\infty = \sup_{x \in K} \left| f(x) - N(x) \right|
\]
So the network can get arbitrarily close to any continuous function, to any desired accuracy.

\end{document}